% ============================================================================
% Chapter 26 --- Wave 8: Workload classification and scale-out
% ----------------------------------------------------------------------------
% Purpose : extend TAC scope to compute and storage workloads at scale.
%           Classification labels become enforced at workload provisioning
%           and at storage binding.
% Page target: ~5 pages.
% Dependencies: Waves 1, 3, 4 required (D17.3); reference layers Ch. 9 + 10;
%               classification scheme institutionally chosen (D10.4); cluster
%               separation by sensitivity (D9.6).
% Mirrors wave template (D18.2).
% ----------------------------------------------------------------------------
% Decisions registered in _coherence-EN.md:
%   D26.1 classification labels enforced at provisioning, retrospective for
%         existing workloads
%   D26.2 cluster separation rule (D9.6) applied during this wave
%   D26.3 snapshots and replicas brought under classification governance (D10.7)
%   D26.4 exit criteria reference Compute + Storage conformance subsets
% ============================================================================

\chapter{Wave~8 --- Workload classification and scale-out}
\label{ch:wave8-workload}

\begin{caveatbox}[Dependencies]
\noindent Wave~8 applies the wave chapter template established in
\trchap{ch:wave0-foundations}{}. Required pre-requisites are
Wave~1 (identity emitting institutional schema), Wave~3 (network
zone catalogue declared and operational), and Wave~4 (policy
plane consuming classification signals); Wave~5 (pilot) strengthens the wave when present. It applies the
contracts of \trchap{ch:compute}{} and \trchap{ch:storage}{},
the cluster separation rule of D9.6, and the
classification-labelling discipline of D10.4 and D10.7. Closure is
gated by the Compute and Storage conformance subsets of
\trchap{ch:conformance}{}.
\end{caveatbox}

Wave~8 is the wave in which classification ceases to be a
documentary discipline and becomes an enforced architectural
property. Workloads --- virtual machines, containers, research
workspaces, administrative services --- carry classification
labels that the policy plane consults at provisioning and that the
storage layer respects at binding. The wave's principal
engineering effort is not in introducing the labels (their
schema was fixed in earlier waves) but in extending classification
to the existing workload estate and in enforcing the cluster-
separation rule that the labels imply.

%-----------------------------------------------------------------------------
\section{Goal and dependencies}
\label{sec:w8-goal}

The goal of Wave~8 is to make the classification of every
workload an authoritative input to its provisioning, its storage
binding, and its policy evaluation. Three concrete objectives
follow:

\begin{itemize}[leftmargin=1.5em,nosep]
  \item every newly provisioned workload declares its
    classification at the moment of provisioning, with the
    declaration verified by the policy layer;
  \item every existing workload acquires a classification through
    a retrospective labelling exercise, with explicit governance
    of the period during which workloads operate without
    classification;
  \item the cluster-separation rule of D9.6 (workloads of
    distinct sensitivity classifications should not share a
    cluster) is operationalised on the production estate, with
    documented exceptions where the rule is impractical.
\end{itemize}

%-----------------------------------------------------------------------------
\section{Pre-requisites}
\label{sec:w8-prereq}

Beyond Waves~1, 3 and 4, four operational conditions strengthen
Wave~8:

\begin{itemize}[leftmargin=1.5em,nosep]
  \item the institutional classification scheme is stable
    (whether TLP-derived, domain-specific, or institutional ---
    the scheme matters less than its stability);
  \item the policy corpus of Wave~4 contains the classification-
    consuming rules that the wave will operationalise;
  \item the storage layer's classification metadata mechanism is
    in place, with a documented procedure for assigning and
    revising labels;
  \item the network zone catalogue of Wave~3 is mature enough to
    accommodate workload migrations between zones when the
    classification of a workload changes.
\end{itemize}

%-----------------------------------------------------------------------------
\section{Coexistence pattern}
\label{sec:w8-coexistence}

Wave~8 introduces enforcement gradually, distinguishing newly-
provisioned workloads from existing ones.

\paragraph{Newly-provisioned workloads.} From the wave's start,
every new workload provisioning request must declare a
classification; the absence of a declaration produces an explicit
provisioning failure that the requester can act on. The
classification is consulted by the policy layer's admission
decision; the storage layer applies the corresponding label to
any volume bound to the workload.

\paragraph{Existing workloads.} A retrospective labelling
exercise classifies the existing workload estate. The exercise is
typically run unit-by-unit or service-by-service, with a
nominated classification owner per cohort. Workloads remain
operational throughout, with an interim default (typically the
most permissive classification consistent with the workload's
visible behaviour) until the formal classification is recorded.

\paragraph{Cluster-separation enforcement.} The rule that
workloads of distinct sensitivity classifications should not
share a cluster (D9.6) is applied as a target rather than as a
hard gate at wave open: existing co-located workloads are
identified, their separation is scheduled, and the schedule is
managed against the institutional change calendar. The wave
closes only when the agreed separation has been achieved or
explicitly excepted.

%-----------------------------------------------------------------------------
\section{Trajectory: greenfield}
\label{sec:w8-greenfield}

For institutions in a greenfield position, Wave~8 establishes
the classification-aware provisioning from the start: there are
no existing workloads to retro-classify, and the cluster-
separation rule applies trivially because clusters are
provisioned per classification class as they are needed.
Indicative duration is four to six months.

%-----------------------------------------------------------------------------
\section{Trajectory: brownfield single-suite-centric}
\label{sec:w8-brownfield-ms}

For institutions whose compute estate is dominated by a single
proprietary vendor suite (public-cloud tenancies aligned to that
vendor, the vendor's productivity and collaboration workloads,
virtualisation substrates with vendor-heavy clusters), Wave~8
brings classification to the visible workload boundary while
accepting that some vendor-internal workloads remain partly
opaque.

\paragraph{Pattern.} The retrospective labelling targets the
workloads visible to the institution's provisioning surface ---
virtual machines under institutional management, containers
under the institution's container-orchestration platform,
research workspaces. The dominant suite's productivity workloads
(mailboxes, document repositories, collaboration workspaces) are
addressed at the application level rather than at the workload
level, with classification expressed as sensitivity labels
integrated with the incumbent identity-and-productivity suite
vendor's information-protection capability where the institution
has adopted it.

\paragraph{Cluster-separation under a dominant-suite estate.}
Virtualisation-based or vendor-private-cloud clusters that host
workloads of mixed classifications have a documented separation
plan. The plan typically aligns with the next infrastructure
refresh cycle or with a new cluster's provisioning, rather than
with disruptive mid-cycle migration.

\paragraph{Indicative duration.} Twelve to twenty-four months,
dominated by the retrospective labelling and by the cluster-
separation schedule.

%-----------------------------------------------------------------------------
\section{Trajectory: brownfield hybrid}
\label{sec:w8-brownfield-hybrid}

For institutions with a multi-vendor compute estate (an
open-source private cloud plus a proprietary virtualisation
platform plus public-cloud burst, with a mix of
container-orchestration clusters across organisational units),
Wave~8's principal effort is the
\emph{normalisation} of classification across these substrates.

\paragraph{Pattern.} A common classification metadata interface
is established across the substrates: the same label vocabulary
is applied whether a workload runs on the open-source private
cloud, on the proprietary virtualisation platform, or
in a cloud burst. The workload-context events of
\trchap{ch:compute}{} carry the classification consistently;
the storage labels of \trchap{ch:storage}{} match. The
multi-vendor reality of the estate is preserved; the
classification governance is unified.

\paragraph{Unit-level classification authorship.} In
research-intensive institutions, classification is often
authored at organisational-unit or research-project level rather
than centrally. The wave operationalises this delegation: a
classification owner per cohort, a central review for
consistency, no central authority over individual labels.

\paragraph{Indicative duration.} Twelve to twenty-four months.

%-----------------------------------------------------------------------------
\section{Reversibility criteria}
\label{sec:w8-reversibility}

Wave~8 is reversible at the per-workload level for the duration
of the retrospective labelling and for an additional grace
period thereafter.

\begin{itemize}[leftmargin=1.5em,nosep]
  \item Per-workload reversal: a classification that proves
    operationally inappropriate is revised through the labelling
    governance rather than through a wave-level reversal.
  \item Cluster-separation reversal: a separation schedule that
    cannot be sustained for resource reasons is documented as an
    exception, not as a wave failure.
  \item Provisioning enforcement toggle: the requirement for a
    classification declaration at provisioning can be disabled
    in defined emergency cases, with the disablement logged and
    reviewed.
\end{itemize}

%-----------------------------------------------------------------------------
\section{Exit criteria}
\label{sec:w8-exit}

Wave~8 closes when the Compute and Storage conformance subsets
of \trchap{ch:conformance}{} (\S\ref{sec:cf-suite}) pass, when
the agreed classification coverage is achieved, and when the
cluster-separation schedule has been completed or explicitly
excepted.

\begin{itemize}[leftmargin=1.5em,nosep]
  \item Compute conformance subset (Kubernetes API conformance,
    OCI image migration, SLSA attestation verification);
    Storage conformance subset (S3 conformance, KMIP migration,
    snapshot inventory, classification preservation across
    export/import).
  \item Classification coverage: the institutionally agreed
    proportion of workloads has a recorded classification.
    Workloads in the residual exception register are documented
    individually with the rationale.
  \item Cluster-separation schedule: closed, with documented
    exceptions where applicable. Exceptions are revisited at
    the next sovereignty grid review.
  \item Snapshots and replicas (D10.7): brought under the
    classification governance, with their location and
    retention reconciled with the primary storage.
\end{itemize}

\begin{realworldbox}[Workload classification at research institutions]
\noindent Multi-tenant research-cloud architectures, such as
those reported under large national research-computing
allocation programmes or within continental-scale federated
research clouds, illustrate
classification-aware patterns at substantial scale. The mapping
of those patterns to a single-institution context is non-trivial
and would benefit from the institution's own pilot work
(Wave~5) rather than from direct transposition.
\end{realworldbox}

%-----------------------------------------------------------------------------
\section{Regulatory anchoring}
\label{sec:w8-regulatory}

\noindent Wave~8 brings classified workloads under the
attestation discipline at scale. It contributes to the
security-of-processing obligation of \loc{data-protection-regime}
(e.g.\ \gls{gdpr}~Art.~32, applied to the workload
tier) and to that regime's data-minimisation principle (e.g.\
\gls{gdpr}~Art.~5(1)(c), realised through classification-aware
tiering). It also operationalises the business-continuity and
secure-development obligations of \loc{cybersecurity-baseline-regime}
(e.g.\ NIS2~Art.~21(2)(c), through declarative
redeployability, and (e), through SLSA-level supply-chain
attestations). Where the institution adopts them, the wave is also
a principal anchor of ISO/IEC~27001~A.8.32 (change
management) at scale and contributes to NIST CSF \textsf{PROTECT}
(PR.PS) and \textsf{RECOVER} (RC.RP). Detailed mapping in
\trchap{ch:regulatory-mapping}{} \S\ref{sec:rm-gdpr},
\S\ref{sec:rm-nis2}, \S\ref{sec:rm-iso},
\S\ref{sec:rm-nist-csf}.
